% s-244-p244-m002; original PDF 244, 245
\Aphorism{s-244-p244-m002}{339}
Der grosse, schöne, schöpferisch überschäumende Mensch, der die Kunst für sich nöthig hat und erfindet, schafft als Begleiterscheinung zur Kunst auch den Inbegriff der Mittel, die ihm für jeden Augenblick, da er dessen bedarf, die lebendige sinnliche Gegenwart des Kunstwerks verbürgen; um des Artistischen willen muss auch ein Apparat für das Technische, Formale, Historische da sein. Dass die Menschen, die hier den Apparat vorstellen, also etwa Musikanten, Ästhetiker, Kritiker, Kunstschulmeister jeder Art, den Gegenwerth für ihre Maschinenarbeit empfangen und von der Kunst leben, ist billig; auch dass sie beim Publicum als die eigentlichen monopolisirten Inhaber von Kunstinteresse und -verständniss gelten, mag hingehen: es ist dies zwar eine Erschleichung und Fiction, wurde aber als solche in den Arbeitsvertrag mit aufgenommen und muss nach aussen hin ebenso vertreten werden, wie etwa von der Monarchie das Gottesgnadenthum, widrigenfalls der ganze Zweck der Institution fraglich wird. Der Künstler ist zu reinlich und verletzlich, um unmittelbar mit dem Publicum zu verkehren; die Einschaltung von „Mittlern“ ist nothwendig, auf die Gefahr hin, dass der Priester dem „Gotte“ die fettesten Bissen des Opfers vom Munde wegstehle. — Aber dass die Kunstprofessionisten und -professoren jene Fiction schliesslich im Ernst aufrechterhalten und sogar ihrem Arbeitgeber ins Angesicht durchsetzen möchten, dass sie \Foreign{in nomine Domini} reden und den ächten Künstler nach ihrem selbstverfassten Exercirreglement drillen wollen — das ist zu komisch und eine so völlige Verkennung der Rangordnung, dass selbst Einer, der den Künstlern nicht jede Anmassung nachsieht, gegen diesen Missbrauch der Licenz Einspruch erheben muss. Der Künstler steht unter Gesetzen, gewiss! nur nicht unter denen seiner eigenen Handlanger; noch ist es nur in der aristophanischen Komödie, nicht in Wirklichkeit dahin gekommen, dass Dionysos sich hinter seinen eignen Priester verkriecht.
\par

